\scanpage{440}
\mspage{522}

\noindent (116)
\diagram{\begin{tikzcd}[ampersand replacement=\&,column sep=small,row sep=small]
 \& \& G_m\arrow[ddrr,dashed,"{\lambda\mapsto\lambda^2}"] \& \& \\
 \& R_{\rho,\sigma}
       \arrow[ur,dashed,"{\kappa}"]
       \arrow[drr,hook,"{PG}"]
 \& \& \& \\
 SR_{\rho,\sigma}
       \arrow[ur,hook,"{G_m}"]
       \arrow[rrr,dashed,"{G}"]
       \arrow[ddrr,hook,"{SG}"']
 \& \& \& \mathcal{M}_{\rho,\sigma}\arrow[r,dashed,"{\chi}"] \& G_m \\
 \& \& \& \& \\
 \& \& \mathcal{M}_{\rho,\sigma}^{+}
       \arrow[uur,hook,"{G_m}"]
 \& G_m\arrow[ddr,dashed,"{\lambda\mapsto\lambda^2}"] \& \\
 \& R_{\rho,\sigma}^{\circ}
       \arrow[uuuu,hook,"{\mu_2\times\mu_2}"]
       \arrow[urr,dashed,"{\kappa^{\circ}}"']
       \arrow[drr,hook,"{PG}"]
 \& \& \& \\
 SR_{\rho,\sigma}^{\circ}
       \arrow[uuuu,hook,"{\mu_2\times\mu_2}"']
       \arrow[ur,hook,"{G_m}"]
       \arrow[rrr,dashed,"{G}"]
       \arrow[ddrr,hook,"{SG}"']
 \& \& \& \mathcal{M}_{\rho,\sigma}^{\circ}
       \arrow[uuuu,hook,"{\mu_2\times\mu_2}"]
       \arrow[r,dashed,"{\chi^{\circ}}"]
 \& G_m \\
 \& \& \& \& \\
 \& \& \mathcal{M}_{\rho,\sigma}^{+\circ}
       \arrow[uuuu,hook,
        "{\begin{gathered}\inserted{\mu_2\times\mu_2}\\
        \deleted{\unsure{\mu_2\times\mu_2}}\end{gathered}}"']
       \arrow[uur,hook,"{G_m}"]
 \& \& \\
 \& \& \rotatebox{90}{$\simeq$} \& \&
\end{tikzcd}}

As in (95)(96), the lower face is deduced from the upper one by
\unsure{passing} to the identity components, the four quotient groups are
isomorphic to one another by the canonical arrows, they are isomorphic to
$\mu_2\times\mu_2$ --- \unsure{itm|item} in the cube \inserted{attached}
below (117). The lateral faces of the cube are cartesian and.
cocartesian, but not the upper faces nor the lower face.

\mspage{439}

\begin{center}
{\small
439\\
\texttt{https://grothendieck.umontpellier.fr}\\
Université de Montpellier, UMR CNRS 5149, Case courrier 051, Place Eugène Bataillon\\
34095 Montpellier cedex 5 - France
}
\end{center}

\scanpage{441}
\mspage{523}

\[
\tag{117}
\diagram{\begin{tikzcd}[ampersand replacement=\&,column sep=1.5em,row sep=2.5em]
 \& \& R_{\rho,\sigma}\arrow[drr,hook,"{{PG}}"] \& \& \& \\
 SR_{\rho,\sigma}\arrow[urr,hook,"{{G_m}}"]\arrow[rrrr,dashed,"{{G}}"]\arrow[dr,hook,"{{SG}}"'] \& \& \& \& \mathcal M_{\rho,\sigma}\arrow[dr,hook] \& \\
 \& \begin{gathered}\mathcal M_{\rho,\sigma}^{+}\\[-.3em]\inserted{\unsure{\simeq SR_{\rho,\sigma}\times SG}}\end{gathered}\arrow[urrr,hook,"{{G_m}}"]\arrow[drr,shift left=.6ex] \& \& \& \& \mathcal T_{\rho,\sigma} \\
 \& \& R_{\rho,\sigma}^{\circ}\arrow[uuu,hook,"{{\mu_2\times\mu_2}}",description,pos=.6]\arrow[drr,hook,"{{SG}}"] \& \mathcal T_{\rho,\sigma}^{+}\arrow[urr,hook,"{{G_m}}"]\arrow[ull,shift left=.6ex] \& \& \\
 \begin{gathered}SR_{\rho,\sigma}^{\circ}\\[-.5em]\rotatebox{90}{$\simeq$}\\[-.5em]ST_\rho\times ST_\sigma\end{gathered}\arrow[uuu,hook,"{{\mu_2\times\mu_2}}"']\arrow[urr,hook,"{{G_m}}"]\arrow[rrrr,dashed,"{{G}}"]\arrow[dr,hook,"{{SG}}"'] \& \& \& \& \mathcal M_{\rho,\sigma}^{\circ}\arrow[uuu,hook,"{{\mu_2\times\mu_2}}",description]\arrow[dr] \& \\
 \& \mathcal M_{\rho,\sigma}^{+\circ}\arrow[uuu,hook,"{{\mu_2\times\mu_2}}"']\arrow[urrr,hook,"{{G_m}}"] \& \& \& \& \mathcal T_{\rho,\sigma}^{\circ}\arrow[uuu,hook,"{{\mu_2\times\mu_2}}"] \\
 \& \& \& \begin{gathered}\mathcal T_{\rho,\sigma}^{+\circ}\\[-.5em]\rotatebox{90}{$\simeq$}\\[-.5em]ST_\rho\times ST_\sigma\end{gathered}\arrow[uuu,hook,"{{\mu_2\times\mu_2}}"]\arrow[ull]\arrow[urr,hook,"{{G_m}}"'] \& \&
\end{tikzcd}}
\]

The composite front left face of the cube is of a particularly simple nature, since the middle terms $\mathcal M_{\rho,\sigma}^{+}$, $\mathcal M_{\rho,\sigma}^{+\circ}$ are identified \unsure{simply} with the direct products of the outer terms
\[
\begin{gathered}
SR_{\rho,\sigma}\xrightarrow{\sim}\mathcal T_{\rho,\sigma}^{+}
\quad\text{and}\quad
SR_{\rho,\sigma}^{\circ}\xrightarrow{\sim}\mathcal T_{\rho,\sigma}^{+\circ}\\
\hphantom{SR_{\rho,\sigma}\xrightarrow{\sim}}\simeq S'T_\rho\times SN'_\sigma
\qquad\qquad\simeq ST_\rho\times ST_\sigma
\end{gathered}
\]
with the group $SG$. This is almost the case too for the parallel face (back right face), \inserted{up to the group $\mu_2$ -- this becomes true \unsure{on}} on condition of replacing the \deleted{groups $\mathcal M_{\rho,\sigma}$} middle terms $\mathcal M_{\rho,\sigma}$, $\mathcal M_{\rho,\sigma}^{\circ}$ by the $\mu_2$-coverings
\[
\tag{118}
\mathcal M_{\rho,\sigma}^{\sim}\simeq\unsure{R}_{\rho,\sigma}\times SG,
\qquad \mathcal M_{\rho,\sigma}^{\circ\sim}\simeq R_{\rho,\sigma}^{\circ}\times SG
\]
(cf. \unsure{(107)|(102)}, (103)), which \unsure{will have} also the virtue

\mspage{440}

\begin{center}
{\small
440\\
\texttt{https://grothendieck.umontpellier.fr}\\
Université de Montpellier, UMR CNRS 5149, Case courrier 051, Place Eugène Bataillon\\
34095 Montpellier cedex 5 - France
}
\end{center}

\scanpage{442}

\unsure{delectable} of making the four squares of the upper and lower faces of the two attached cubes cartesian and cocartesian, which they were not before, \unsure{regrettably}.

To finish these scales on the \unsure{radicals} I am going to introduce a canonical homomorphism
\[
\tag{118}
\diagram{\begin{tikzcd}[ampersand replacement=\&]
G_m\arrow[r,hook,"{{i_R}}"] \& R_{\rho,\sigma}^{\circ}
\end{tikzcd}}
\]
via a canonical homomorphism
\[
\tag{119}
\diagram{\begin{tikzcd}[ampersand replacement=\&,row sep=1em]
G_m\arrow[r,hook,"{{i_{\mathcal T}}}"] \& \mathcal T_{\rho,\sigma}^{\circ}\simeq T_\rho\times_{G_m}T_\sigma \\
\lambda\arrow[r,mapsto] \& (\lambda,\lambda)
\end{tikzcd}}
\]
of which \unsure{the composite|the composite} with
\[
\diagram{\begin{tikzcd}[ampersand replacement=\&]
\chi:\mathcal T_{\rho,\sigma}^{\circ}\arrow[r] \& G_m
\end{tikzcd}}
\]
is indeed given by $\lambda\mapsto\lambda^2$ \inserted{(104)}, so that thanks to the description (105) of $R_{\rho,\sigma}$ \inserted{in terms of $\mathcal T_{\rho,\sigma}$} as a fiber product (\unsure{paraphrased} as an expression of $R_{\rho,\sigma}^{\circ}$ in terms of $\mathcal T_{\rho,\sigma}^{\circ}$) we have a unique homomorphism (118), making the diagram commutative

\mspage{441}

\begin{center}
{\small
441\\
\texttt{https://grothendieck.umontpellier.fr}\\
Université de Montpellier, UMR CNRS 5149, Case courrier 051, Place Eugène Bataillon\\
34095 Montpellier cedex 5 - France
}
\end{center}

\scanpage{443}
\mspage{524}

\[
\tag{120}
\diagram{\begin{tikzcd}[ampersand replacement=\&,column sep=3em,row sep=2em]
G_m\arrow[r,"{{i_{\mathcal T}}}"]\arrow[dr,"{{i_R}}"'] \& \mathcal T_{\rho,\sigma}^{\circ} \\
\& R_{\rho,\sigma}^{\circ}\arrow[u,"{{\mathrm{can}}}"']
\end{tikzcd}}
\]
or, more prettily, the diagram of exact sequences
\[
\tag{121}
\diagram{\begin{tikzcd}[ampersand replacement=\&,column sep=1.8em,row sep=2em]
1\arrow[r] \& \mu_2\arrow[r] \& G_m\arrow[r,"{{\lambda\mapsto\lambda^2}}"]\arrow[d,no head] \& G_m\arrow[r,no head] \& 1 \\
\& \& {} \& \&
\end{tikzcd}}.
\]

By construction, we find at once
\[
\tag{122}
\diagram{\begin{tikzcd}[ampersand replacement=\&,column sep=3em]
G_m\arrow[r,"{{i_R}}"]\arrow[rr,bend right=28,"{{\mathrm{id}}}"'] \& R_{\rho,\sigma}^{\circ}\arrow[r,"{{\kappa}}"] \& G_m
\end{tikzcd}}
\]
whence direct product decompositions
\[
\tag{123}
\left\{\begin{aligned}
R_{\rho,\sigma}&\simeq SR_{\rho,\sigma}\times G_m\simeq S'T_\rho\times SN'_\sigma\times G_m\\
R_{\rho,\sigma}^{\circ}&\simeq SR_{\rho,\sigma}^{\circ}\times G_m\\
&\simeq ST_\rho\times ST_\sigma\times G_m
\end{aligned}\right.,
\]
which considerably sharpens (96) with which I was nevertheless so pleased!

In terms of this decomposition, we are going to make explicit the \deleted{\unsure{h}}homomorphism
\[
\diagram{\begin{tikzcd}[ampersand replacement=\&,column sep=1.7em]
\mu_2\arrow[r,hook] \& R_{\rho,\sigma}^{\circ}\subset R_{\rho,\sigma}
\end{tikzcd}}
\quad
\left(\text{\unsure{and }}
\diagram{\begin{tikzcd}[ampersand replacement=\&,column sep=1.4em]
R_{\rho,\sigma}\arrow[r,no head] \& \mathcal T_{\rho,\sigma}
\end{tikzcd}},\ 
\diagram{\begin{tikzcd}[ampersand replacement=\&,column sep=1.4em]
R_{\rho,\sigma}^{\circ}\arrow[r] \& \mathcal T_{\rho,\sigma}^{\circ}
\end{tikzcd}}\right)
\]
which occurs in (102), (103), $\mu_2=R_{\rho,\sigma}\cap SG=R_{\rho,\sigma}^{\circ}\cap SG$,

\deleted{\unsure{We find} in the explicit description (104), that $\mu_2$ corresponds to the $\varepsilon\in\Gamma\mu_2$, \unsure{that} $\varepsilon$ \illegible{2}}

\mspage{442}

\begin{center}
{\small
442\\
\texttt{https://grothendieck.umontpellier.fr}\\
Université de Montpellier, UMR CNRS 5149, Case courrier 051, Place Eugène Bataillon\\
34095 Montpellier cedex 5 - France
}
\end{center}

\scanpage{444}

For the homomorphism
\[
\diagram{\begin{tikzcd}[ampersand replacement=\&,column sep=2em]
R_{\rho,\sigma}\arrow[r] \&
\begin{gathered}
\inserted{\{(\underline a,\underline b)\mid(\det\underline a/\det\underline b)^2=1\}}\\[-.5em]
\rotatebox{90}{$=$}\\[-.5em]
T_{\rho,\sigma}
\end{gathered}\arrow[r,hook] \& T_\rho\times N'_\sigma
\end{tikzcd}}
\]
it suffices to \deleted{define} \inserted{make explicit} the composites
\[
\tag{124}
\left\{\begin{gathered}
\diagram{\begin{tikzcd}[ampersand replacement=\&,column sep=2.2em,row sep=1em]
\& \& T_\rho \\
R_{\rho,\sigma}\arrow[r]\arrow[d,phantom,"{{\rotatebox{90}{$\simeq$}}}" description] \& \unsure{T_{\sigma\sigma}|T_{\rho,\sigma}}\arrow[ur]\arrow[dr] \& \\
S'T_\rho\times SN'_\sigma\times G_m \& \& N'_\sigma
\end{tikzcd}}\\[1em]
\diagram{\begin{tikzcd}[ampersand replacement=\&]
(\underline a,\underline b,\lambda)\arrow[r] \& (\lambda\underline a,\lambda\underline b)
\end{tikzcd}}
\end{gathered}\right.
\]
\margin{we have as well
\[
\underbrace{\det(\lambda\underline a)^2}_{\substack{\lambda^2\\\text{since }(\det\underline a)^2=1}}
=
\underbrace{\det(\lambda\underline b)^2}_{\substack{\lambda^2\ \text{since}\\\det\underline b=1}}
\]}

The kernel of this homomorphism therefore consists of the $(\unsure{\lambda}^{-1},\lambda^{-1},\lambda)$, \unsure{with|where} $\lambda\in\Gamma G_m$ such that $\lambda^{-1}\in\Gamma S'T_\rho$, $\lambda^{-1}\in\Gamma SN'_\sigma$, \deleted{\unsure{that is|and}} i.e. $(\lambda^2)^2=1$ and $\lambda^2=1$, i.e. $\lambda\in\Gamma\mu_2$. Thus
\[
\tag{125}
\left\{\begin{gathered}
\diagram{\begin{tikzcd}[ampersand replacement=\&,column sep=2em]
\begin{gathered}\mu_2\\[-.5em]\rotatebox{90}{$=$}\\[-.5em]R_{\rho,\sigma}\cap SG\end{gathered}\arrow[r,"{{\mathrm{can}}}"] \&
\underbrace{R_{\rho,\sigma}^{\circ}}_{ST_\rho\times ST_\sigma\times G_m}\subset
\underbrace{R_{\rho,\sigma}}_{S'T_\rho\times SN'_\sigma\times G_m}
\end{tikzcd}}\\[1em]
\diagram{\begin{tikzcd}[ampersand replacement=\&]
\lambda\arrow[r,mapsto] \& (\lambda,\lambda,\lambda)
\end{tikzcd}}
\end{gathered}\right.\,.
\]

We thus find the particularly simple explicit description
\[
\tag{126}
\diagram{\begin{tikzcd}[ampersand replacement=\&,column sep=0em,row sep=.8em]
\mathcal T_{\rho,\sigma}\simeq(S'T_\rho\times SN'_\sigma\times G_m)/ \& \mu_2 \\
\& \begin{gathered}\mu_2\ \text{sent ``diagonally'' into the}\\\text{triple product}\end{gathered}\arrow[u]
\end{tikzcd}}
\]

We will note that we have canonical homomorphisms
\[
\diagram{\begin{tikzcd}[ampersand replacement=\&,column sep=1.2em]
\mu_2\arrow[r] \& ST_\rho\subset S'T_\rho
\end{tikzcd}},\quad
\diagram{\begin{tikzcd}[ampersand replacement=\&,column sep=1.2em]
\mu_2\arrow[r] \& ST_\sigma\subset SN'_\sigma
\end{tikzcd}},\quad
\diagram{\begin{tikzcd}[ampersand replacement=\&,column sep=1.2em]
\mu_2\arrow[r,no head] \& G_m
\end{tikzcd}}
\]
(which are used in (125)) whence
\[
\tag{127}
\left\{\begin{gathered}
\diagram{\begin{tikzcd}[ampersand replacement=\&,column sep=1.5em]
\mu_2\times\mu_2\times\mu_2\arrow[r] \&
\underbrace{ST_\rho\times ST_\sigma\times G_m}_{R_{\rho,\sigma}^{\circ}}\subset
\underbrace{S'T_\rho\times SN'_\sigma\times G_m}_{R_{\rho,\sigma}}
\end{tikzcd}}\\
\diagram{\begin{tikzcd}[ampersand replacement=\&]
(\lambda',\lambda'',\lambda)\arrow[r,mapsto] \& (\lambda',\lambda'',\lambda)
\end{tikzcd}}
\end{gathered}\right.
\]
which is none other than (48))

\mspage{443}

\begin{center}
{\small
443\\
\texttt{https://grothendieck.umontpellier.fr}\\
Université de Montpellier, UMR CNRS 5149, Case courrier 051, Place Eugène Bataillon\\
34095 Montpellier cedex 5 - France
}
\end{center}